% =====================================================================
%  CHAPTER 12: 食品安全性风险分析和控制
% =====================================================================
\chapter{食品安全性风险分析和控制}
\setcounter{chapter}{12}
\chapterintro{
\keyword{掌握}风险分析框架（危害识别、危害特征描述、暴露评估、风险特征描述）；掌握NOEL/NOAEL/ADI。\keyword{熟悉}营养毒理学。
}

% =====================================================================
%  第一节 食品安全性风险分析框架
% =====================================================================
\section{食品安全性风险分析框架}

\subsection{风险分析框架的四大步骤}

\begin{keybox}
\textbf{食品安全性风险分析框架（四大步骤*）：}

\begin{enumerate}[leftmargin=*]
    \item \textbf{危害识别（Hazard Identification）：}确定人体摄入某种食品中的危害因素与健康不良效应之间是否存在因果关系。信息来源：流行病学研究、动物毒理学试验、体外试验、构效关系分析。

    \item \textbf{危害特征描述（Hazard Characterization）：}对危害因素可能引起的不良健康效应进行定性和/或定量评价。核心任务：建立剂量-反应关系（Dose-Response Relationship），确定NOAEL或BMDL。

    \item \textbf{暴露评估（Exposure Assessment）：}对通过食品或其他途径摄入人体的生物、化学和物理因素的量进行评估。膳食暴露量 = $\sum$(食品中化学物浓度 $\times$ 食品消费量) / 体重。

    \item \textbf{风险特征描述（Risk Characterization）：}综合前三步的结果，分析和评估危害因素对人群健康产生不良效应的风险和严重程度。与健康指导值比较：暴露量 $<$ HBGV $\to$ 风险可接受。
\end{enumerate}
\end{keybox}

% =====================================================================
%  第二节 营养毒理学
% =====================================================================
\section{营养毒理学}

\begin{defbox}
\textbf{营养毒理学（Nutritional Toxicology）：}研究营养素过量摄入的毒副作用。UL（可耐受最高摄入量）是营养毒理学的重要指标。

\textbf{三大研究方向：}
\begin{enumerate}[leftmargin=*]
    \item 营养素过量摄入的不良反应及UL的制定。
    \item 营养素对外源化学物代谢和毒性的影响。
    \item 膳食中有毒物质对营养素代谢和营养过程的影响。
\end{enumerate}
\end{defbox}

% =====================================================================
%  第三节 安全性评价方法
% =====================================================================
\section{安全性评价方法}

\begin{defbox}
\textbf{安全性评价方法：}

\begin{enumerate}[leftmargin=*]
    \item \textbf{急性毒性试验（LD$_{50}$）：}了解受试物的急性毒性强度和性质，测定半数致死量。
    \item \textbf{30天喂养试验：}了解多次重复染毒条件下受试物的毒性。
    \item \textbf{90天喂养试验：}观察较长时间重复染毒的毒性效应，确定NOAEL。
    \item \textbf{传统致畸试验：}评价受试物是否具有致畸作用。
    \item \textbf{慢性毒性试验（包括致癌试验）：}观察长期（终生）染毒的慢性毒性和致癌性。
\end{enumerate}
\end{defbox}

% =====================================================================
%  第四节 关键毒理学指标
% =====================================================================
\section{关键毒理学指标}

\subsection{最大无作用剂量}

\begin{keybox}
\textbf{最大无作用剂量（MNL/NOEL/NOAEL）*：}在一定时间内，某种外源化学物按一定方式与机体接触，用现代检测方法和最灵敏的观察指标不能检出任何有害作用的最高剂量。

\begin{itemize}[leftmargin=*]
    \item \textbf{MNL（Maximum No-effect Level）：}最大无作用剂量。
    \item \textbf{NOEL（No Observed Effect Level）：}未观察到作用水平。
    \item \textbf{NOAEL（No Observed Adverse Effect Level）：}未观察到有害作用水平。
\end{itemize}
\end{keybox}

\subsection{每日允许摄入量}

\begin{keybox}
\textbf{每日允许摄入量（Acceptable Daily Intake, ADI）*：}人类每日摄入某物质直至终生，根据当时已知事实，不会产生可察觉健康危害的量。单位：mg/(kg bw$\cdot$d)。

\textbf{推导公式：}
\[
\text{ADI} = \frac{\text{NOAEL}}{\text{UFs}}
\]

其中UFs（Uncertainty Factors，不确定系数）常规取100倍（种间差异10 $\times$ 个体差异10）。
\end{keybox}

% =====================================================================
%  复习题
% =====================================================================
\reviewcontent{
\section{复习题}

\subsection*{一、名词解释}

\begin{enumerate}[leftmargin=*, itemsep=6pt]
    \item \textbf{危害识别（Hazard Identification）、危害特征描述（Hazard Characterization）}
    \item \textbf{暴露评估（Exposure Assessment）、风险特征描述（Risk Characterization）}
    \item \textbf{营养毒理学（Nutritional Toxicology）}
    \item \textbf{MNL、NOEL、NOAEL}
    \item \textbf{ADI（每日允许摄入量）}
    \item \textbf{UFs（不确定系数）}
\end{enumerate}

\subsection*{二、简答题}

\begin{enumerate}[leftmargin=*, itemsep=8pt]
    \item \textbf{【重点】}简述食品安全性风险分析框架的四大步骤及各步骤的主要任务。
    \item \textbf{【重点】}简述MNL/NOEL/NOAEL的定义及ADI的定义和推导公式。
    \item 简述营养毒理学的概念和主要研究内容。
    \item 简述食品安全性毒理学评价的主要方法。
\end{enumerate}

\subsection*{参考答案要点}

\subsubsection*{名词解释（要点）}

\begin{enumerate}[leftmargin=*, itemsep=4pt]
    \item \textbf{危害识别：}确定危害因素与健康不良效应之间是否存在因果关系。\textbf{危害特征描述：}对不良健康效应进行定性和/或定量评价，建立剂量-反应关系。
    \item \textbf{暴露评估：}对通过食品摄入人体的因素的量进行评估。\textbf{风险特征描述：}综合前三步，评估危害因素对人群健康产生不良效应的风险和严重程度。
    \item \textbf{营养毒理学：}研究营养素过量摄入的毒副作用。UL是营养毒理学的重要指标。
    \item \textbf{MNL/NOEL/NOAEL：}在一定时间内，某种外源化学物按一定方式与机体接触，用现代检测方法和最灵敏的观察指标不能检出任何有害作用的最高剂量。
    \item \textbf{ADI：}人类每日摄入某物质直至终生，根据当时已知事实，不会产生可察觉健康危害的量。ADI = NOAEL / UFs。
    \item \textbf{UFs：}不确定系数，常规取100倍（种间差异10 $\times$ 个体差异10）。
\end{enumerate}

\subsubsection*{简答题（要点）}

\begin{enumerate}[leftmargin=*, itemsep=6pt]
    \item \textbf{四大步骤：}①危害识别（确定因果关系）；②危害特征描述（定性和/或定量评价，建立剂量-反应关系）；③暴露评估（评估摄入量）；④风险特征描述（综合评估风险和严重程度）。
    \item \textbf{MNL/NOEL/NOAEL：}不能检出任何有害作用的最高剂量。\textbf{ADI：}人类每日摄入某物质直至终生不会产生可察觉健康危害的量。\textbf{公式：}ADI = NOAEL / UFs（常规100倍）。
    \item \textbf{营养毒理学：}研究营养素过量摄入的毒副作用。三大方向：①营养素过量不良反应及UL制定；②营养素对外源化学物代谢和毒性的影响；③膳食有毒物质对营养素代谢的影响。
    \item \textbf{方法：}急性毒性试验（LD$_{50}$）、30天喂养试验、90天喂养试验、传统致畸试验、慢性毒性试验（包括致癌试验）。
\end{enumerate}

\newpage
}
